% ---------------------------------------------------------------------------
%  Report template
%  "Seeing Through Solid Rock" -- MYSSP Project P-1
%
%  Compile with pdflatex (twice, so that references resolve).
%  Target length: 6-8 pages including figures. Shorter and sharper beats
%  longer and padded -- the Khufu discovery paper itself is 5 pages.
% ---------------------------------------------------------------------------
\documentclass[11pt,a4paper]{article}
\usepackage[margin=2.5cm]{geometry}
\usepackage{graphicx}
\usepackage{amsmath}
\usepackage{siunitx}
\usepackage[hidelinks]{hyperref}
\usepackage{caption}

\title{Detecting a Hidden Chamber with Cosmic-Ray Muons:\\
       A Simulation Study}
\author{Your Name \\ \small MYSSP Project P-1 \\
        \small Mentor: Assoc.\ Prof.\ Dr.\ Khaw Kim Siang, Tsung-Dao Lee Institute, SJTU}
\date{\today}

\begin{document}
\maketitle

\begin{abstract}
% Four sentences, written LAST.
% 1. What problem. 2. What you did. 3. What number you got. 4. What it means.
% Example shape:
% "Muography images dense objects by counting cosmic-ray muons that pass
%  through them. We simulated ... We find that a detector of area X requires
%  Y days for the expected signal from a Z-metre-radius chamber to reach 5 sigma under the project model. State the chamber radius explicitly."
\end{abstract}

\section{Introduction}
% Why can muons see through rock, and why is that useful?
% One paragraph of physics, one paragraph of applications, one sentence
% stating the question this report answers. Cite Khufu here.

\section{Method}

\subsection{Muon flux and energy loss}
% The Gaisser parametrisation and the a+bE energy-loss model.
% State the formulae, define every symbol, give the units.
\begin{equation}
    E_{\min}(X) = \frac{a}{b}\left( e^{bX} - 1 \right)
\end{equation}
% ... and explain what a, b and X are.

\subsection{Target geometry and ray tracing}
% Describe the pyramid, the void, the detector position.
% Include your density-slice figure here -- it does more than a paragraph.

\subsection{Counting statistics}
% How expected counts are computed. Define expected significance and observed significance, and explain why the headline exposure time uses the expected-significance curve.

\section{Validation}
% This section is what separates a report from a description.
% What did you check, and against what? Step-size convergence, the
% no-void control, hand calculations. Say what agreed and to what precision.

\section{Results}

\subsection{Opacity and count maps}
\begin{figure}[htbp]
    \centering
    % \includegraphics[width=0.9\textwidth]{figures/your_figure.pdf}
    \caption{Write captions that can be understood on their own, without the
             main text. A reader who looks only at your figures should still
             learn your result.}
    \label{fig:maps}
\end{figure}

\subsection{Exposure time required}
% THE headline result. Report the first crossing of 5 sigma on the expected-significance curve. Explain the interpolation/selection method you used; do not invent an uncertainty on the exposure time unless your mentor gives you a defined method.

\subsection{Parameter studies}
% Your two chosen studies. Prediction first, then result, then explanation.

\section{Discussion}
% What limits the measurement? What did the simulation leave out
% (scattering, backgrounds, detector efficiency, uncertainty in rock density)?
% Which of those would matter most in a real campaign, and why?
% Be specific and be honest. This section is read most carefully.

\section{Conclusion}
% Three sentences. No new information.

\section*{Code availability}
% Point to your repository. Say which notebook produces which figure.

\begin{thebibliography}{9}
\bibitem{khufu} K.~Morishima \textit{et al.},
  ``Discovery of a big void in Khufu's Pyramid by observation of cosmic-ray
  muons,'' \textit{Nature} \textbf{552}, 386--390 (2017).
\bibitem{bonomi} G.~Bonomi \textit{et al.},
  ``Applications of cosmic-ray muons,''
  \textit{Prog.\ Part.\ Nucl.\ Phys.} \textbf{112}, 103768 (2020).
\bibitem{pdg} Particle Data Group, ``Cosmic Rays,''
  \textit{Review of Particle Physics}.
\end{thebibliography}

\end{document}
